\documentclass[sigconf,nonacm]{acmart}

%% ---- Packages -------------------------------------------------------
\usepackage{booktabs}
\usepackage{multirow}
\usepackage{amsmath}
\usepackage{graphicx}
\usepackage{xcolor}
\usepackage{hyperref}
\usepackage{cleveref}
\usepackage{microtype}
\usepackage{algorithm}
\usepackage{algorithmic}
\usepackage{subcaption}

%% ---- Metadata -------------------------------------------------------
\title{VINNA: Verifiable Intelligence for Navigating and Analyzing\\
       Construction Sites from Egocentric Video}

\author{VINNA Team}
\affiliation{%
  \institution{HackTech 2025 --- Caltech}
  \city{Pasadena}
  \state{California}
  \country{USA}
}
\email{vinna@hacktech.io}

\begin{document}

\begin{abstract}
Construction sites are spatially complex, dynamic environments where
workers, heavy equipment, materials, and structural elements continuously
reconfigure across three dimensions.
Understanding this spatial evolution---distances between entities,
zone occupancy patterns, temporal changes in site layout---is a
prerequisite for every downstream task in construction intelligence,
from safety compliance to progress monitoring to resource allocation.
Yet current vision systems for construction treat spatial reasoning
as an afterthought, jumping directly to narrow classification tasks
(PPE detection, violation flagging) without building a general spatial
representation.

We introduce VINNA (\textbf{V}erifiable \textbf{I}ntelligence for
\textbf{N}avigating and \textbf{A}nalyzing construction sites), a
framework for egocentric spatial reasoning on construction footage.
VINNA extracts metric distance estimates, spatial relationship graphs,
zone activity classifications, and temporal change signals from
body-worn camera and LiDAR streams.
Crucially, each spatial claim is paired with a \emph{verifiable reward}:
distance estimates are checked against fused depth data, spatial
relationships are validated by 3D geometry, and change detection outputs
are scored against frame-pair differencing.
This verifiability enables self-supervised training via GRPO without
human annotation.

On 638 frames (21 minutes) of Ironsite egocentric fisheye + LiDAR
footage from an active masonry site, organized into 39 temporal batches
with 3072-dimensional Gemini Embedding~2 representations, VINNA produces
spatially grounded scene analyses that cluster into semantically
meaningful groups (within-cluster cosine similarities of 0.4246, 0.4127,
and 0.3765 across three spatial condition tiers).
As a downstream demonstration, we show that VINNA's spatial
representations enable CII compliance classification (P/C/NC) at
84\% accuracy on a 25-frame annotated subset---without any
compliance-specific training.

This paper is an honest report of a 48-hour hackathon prototype with
real construction data, not a polished conference submission.
We describe what works, what does not, and what the spatial intelligence
gap in construction AI looks like from the ground up.
\end{abstract}

%%
%% Keywords
%%
\keywords{spatial intelligence, egocentric video, construction,
          distance estimation, change detection, verifiable rewards,
          VLM, GRPO, LiDAR, spatial reasoning}

\maketitle

%% =====================================================================
\section{Introduction}
\label{sec:intro}

Construction is a \$13~trillion global industry where the physical
environment changes daily~\cite{golparvar2011monitoring}.
Every decision on a construction site---where to place materials,
whether a scaffold is safe, when to pour concrete, how to route
equipment---is fundamentally a \emph{spatial} decision.
Yet the computer vision systems deployed on construction sites today
largely bypass spatial reasoning entirely.

The dominant paradigm is narrow classification: detect hard hats,
flag missing guardrails, count workers in a
zone~\cite{nath2020ppe,yang2018worker}.
These systems answer ``what is present?''\ but not ``where is it
relative to everything else, and how has that changed?''
The spatial questions---\emph{How far is the worker from the edge?
Has the scaffold been moved since yesterday?
Which zone has the highest activity density this hour?}---are the
ones that actually drive operational decisions.

This gap between spatial questions and classification answers is not
unique to construction.
The broader vision-language model (VLM) community has recently
recognized \emph{spatial intelligence} as a first-class research
problem.
SpatialVLM~\cite{spatialvlm2024} demonstrated that VLMs can learn
metric distance estimation from synthetic spatial QA data generated
via monocular depth.
SpatialRGPT~\cite{spatialrgpt2024} extended this to region-grounded
spatial queries over 3D scene graphs.
Ego4D~\cite{ego4d2022} established egocentric video understanding as
a research domain, though its construction coverage is limited to
activity recognition rather than spatial analysis.

None of these systems target construction, and none address the
specific spatial challenges of the domain: extreme fisheye distortion,
LiDAR occlusion from equipment, rapid scene change from active work,
and the need for \emph{metric} (not just relative) spatial claims
grounded in site coordinates.

Meanwhile, the reinforcement learning community has converged on
\emph{verifiable rewards} as a superior alternative to human preference
labels for training language models on tasks with checkable
answers~\cite{faithfulgrpo2026}.
Spatial reasoning is precisely such a task: if a model claims ``the
worker is 2.3 meters from the scaffold edge,'' we can verify that
claim against the fused LiDAR depth data.
Smooth Operator~\cite{smoothoperator2025} showed that continuous
sigmoid-based rewards (SNRA) provide stable training gradients for
numerical spatial predictions---exactly the reward shape we need for
distance estimation.

\paragraph{The VINNA thesis.}
Construction sites need \emph{spatial intelligence}, not violation
classifiers.
Spatial understanding---distance estimation, change detection, spatial
relationship mapping, zone activity analysis---is the foundation.
Compliance classification, progress monitoring, resource optimization,
and safety alerting are all \emph{downstream applications} of that
foundation.
A system that understands space can do all of these; a system that
only detects hard hats can do one.

\paragraph{Contributions.}
\begin{enumerate}
  \item We frame construction-site computer vision as a \emph{spatial
        intelligence} problem and identify four core spatial
        capabilities: metric distance estimation, spatial relationship
        extraction, temporal change detection, and zone activity
        analysis.
  \item We design verifiable reward functions for each spatial
        capability---distance rewards checked against LiDAR depth,
        relationship rewards validated by 3D geometry, change detection
        rewards scored by frame-pair differencing---enabling
        self-supervised GRPO training without human labels.
  \item We demonstrate the approach on 638 frames of real Ironsite
        egocentric footage, showing that spatial representations
        produce meaningful embedding-space structure (39 temporal
        batches, 3072-dim embeddings).
  \item We show CII compliance classification as a downstream task
        that emerges from spatial understanding, achieving 84\% accuracy
        without compliance-specific training.
\end{enumerate}

\paragraph{Scope and honesty.}
VINNA is a 48-hour hackathon prototype.
We have real data, real embeddings, and a real pipeline.
We do not have a trained RL model---the GRPO training loop is designed
but not yet executed.
We report what we built, what the data shows, and what the spatial
intelligence framing makes possible.

%% =====================================================================
\section{Related Work}
\label{sec:related}

\subsection{Spatial Reasoning in Vision--Language Models}

The spatial intelligence gap in VLMs has attracted significant recent
attention.

\emph{SpatialVLM}~\cite{spatialvlm2024} generates 2 billion spatial QA
pairs from monocular depth estimation, training VLMs to answer metric
distance and spatial relationship questions.
Their data generation pipeline---run depth estimation on images, extract
spatial facts from the depth map, template into QA pairs---is directly
applicable to construction footage.
Our work extends their insight to the egocentric fisheye domain with
LiDAR-fused depth (rather than monocular estimates) as ground truth.

\emph{SpatialRGPT}~\cite{spatialrgpt2024} adds region-conditioned
grounding to spatial VLM reasoning: given a bounding box, the model
answers spatial queries about that specific region.
This capability is essential for construction, where the relevant
spatial question is not ``how big is the scene?''\ but ``how far is
\emph{this worker} from \emph{that edge}?''
SpatialRGPT's depth-plugin architecture (inject depth features into a
frozen VLM) informs our sensor fusion strategy.

\emph{ViGoRL}~\cite{vigorl2025} introduces visually grounded
reinforcement learning, requiring each reasoning step to point at a
spatial region in the image.
This grounding constraint prevents the model from generating
spatially unmoored claims---a critical failure mode in our domain,
where a claim like ``worker is 1.5m from the edge'' must correspond
to actual visual evidence.

\subsection{Verifiable Rewards and GRPO}

\emph{Faithful GRPO}~\cite{faithfulgrpo2026} restricts RL training to
tasks with deterministically verifiable rewards, demonstrating that
verifiable correctness is a stronger signal than human preference.
They report a reduction in reasoning inconsistency from 24.5\% to 1.7\%
on Qwen2.5-VL-7B.
Our spatial rewards (distance estimation, relationship classification)
are precisely the kind of verifiable signal Faithful GRPO targets:
there is a ground-truth answer computable from the depth data, and the
model's claim either matches it or does not.

\emph{Smooth Operator}~\cite{smoothoperator2025} addresses the
challenge of training on continuous spatial quantities (distances, areas,
angles) where binary right/wrong rewards provide a sparse gradient.
Their SNRA (Smooth Numerical Reward Approximation) applies a sigmoid
to the absolute error: $R = \sigma(1 - |predicted - actual| / \theta)$,
producing a dense, continuous reward signal that enables stable GRPO
training.
They validate on their Numerical3D-50k dataset of 3D spatial reasoning
tasks.
We adopt the SNRA formulation for distance estimation rewards and extend
it to construction-specific spatial quantities.

\subsection{Egocentric Video Understanding}

\emph{Ego4D}~\cite{ego4d2022} provides 3,670 hours of egocentric video
spanning daily activities, though construction-specific content is
limited.
Their task taxonomy---episodic memory (``when did I last see X?''),
hand-object interaction, social interaction, forecasting---partially
overlaps with construction spatial reasoning.
The episodic memory task, in particular, maps to construction change
detection: ``what changed in zone~B since the morning walkthrough?''

Our egocentric setting inherits the core challenges identified by Ego4D:
extreme viewpoint variation as the wearer moves, frequent occlusion from
the wearer's own body and held objects, and the absence of a stable
camera frame.
We add construction-specific challenges: fisheye lens distortion (180$^\circ$
FOV), dense clutter from materials and equipment, and the need for
metric (not pixel-space) spatial claims.

\subsection{Construction Site Computer Vision}

Prior work in construction CV has focused on progress
monitoring~\cite{golparvar2011monitoring}, worker
detection~\cite{yang2018worker}, and PPE
classification~\cite{nath2020ppe}.
These are important applications, but they share a common limitation:
they jump directly from pixels to a task-specific classifier without
building a general spatial representation of the site.

\emph{Safe-Construct}~\cite{safeconstruct2025} is the closest prior
work, performing 3D multi-view construction violation recognition.
However, Safe-Construct frames the problem as violation detection---the
output is a set of safety violations, not a spatial understanding of the
site.
VINNA inverts this: spatial understanding is the core output, and
violation detection is one of many possible downstream applications.
A system that understands ``the worker is 0.9\,m from an unguarded edge
at 4.2\,m elevation'' can detect the fall hazard, but it can also
estimate scaffold load proximity, plan equipment routing, and track
zone occupancy---none of which a violation classifier can do.

%% =====================================================================
\section{Method}
\label{sec:method}

\subsection{Spatial Intelligence Capabilities}
\label{sec:capabilities}

We decompose construction spatial intelligence into four core
capabilities, each with a corresponding verifiable reward signal.

\subsubsection{Metric Distance Estimation}

Given an egocentric frame and fused depth data, estimate the metric
distance (in meters) between specified entity pairs: worker--edge,
worker--equipment, material--zone boundary, equipment--structure.

This is the most fundamental spatial capability.
It requires the model to (a) identify the relevant entities,
(b) localize them in 3D, and (c) compute the metric separation.
The verifiable reward is direct: compare the predicted distance against
the fused LiDAR measurement.

\subsubsection{Spatial Relationship Extraction}

Extract a spatial relationship graph from each frame: which entities
are above/below, left/right, in front/behind, inside/outside, and
adjacent to each other.
Relationships are typed (``worker \emph{on} scaffold'', ``materials
\emph{near} edge'', ``equipment \emph{inside} exclusion zone'').

This capability builds a structured representation of the scene that
supports compositional spatial queries (``which workers are on
scaffolds near unguarded edges?'').

\subsubsection{Temporal Change Detection}

Given a pair of frames from different times (or different walkthroughs),
identify what has changed: new objects, removed objects, moved objects,
and changed spatial relationships.

This capability is critical for construction, where the site changes
daily.
Progress monitoring, safety auditing, and resource tracking all reduce
to change detection over spatial representations.

\subsubsection{Zone Activity Analysis}

Partition the visible scene into functional zones (work area, storage,
transit corridor, hazard zone) and classify the activity level and
type in each zone.

This capability supports site management decisions: which zones are
active, which are idle, where are workers concentrated, is equipment
blocking transit routes.

\subsection{Verifiable Spatial Rewards}
\label{sec:rewards}

The key insight enabling self-supervised training is that each spatial
capability produces claims that can be verified against geometric ground
truth.
We design three reward functions, one for each verifiable spatial
dimension.

\subsubsection{Distance Estimation Reward (SNRA)}

Following Smooth Operator~\cite{smoothoperator2025}, we use a smooth
sigmoid reward for distance estimation:

\begin{equation}
  R_{\text{dist}}(\hat{d}, d^*) =
    \sigma\!\left(1 - \frac{|\hat{d} - d^*|}{\theta}\right)
  \label{eq:dist_reward}
\end{equation}

where $\hat{d}$ is the model's predicted distance, $d^*$ is the
LiDAR ground-truth distance, $\theta$ is a task-specific tolerance
(e.g., $\theta = 0.5$\,m for worker--edge proximity, $\theta = 2.0$\,m
for equipment--structure separation), and $\sigma$ is the logistic
sigmoid.

This produces a continuous reward in $(0, 1)$ that is maximal when
$\hat{d} = d^*$ and degrades smoothly with increasing error.
Unlike a binary ``correct within threshold'' reward, SNRA provides a
gradient everywhere, enabling stable GRPO optimization.

For fine-grained proximity estimation (e.g., worker distance to an
unprotected edge), we additionally define a binary boundary reward:

\begin{equation}
  R_{\text{boundary}}(\hat{d}, d^*; \tau) =
    \mathbf{1}\!\left[(\hat{d} < \tau) = (d^* < \tau)\right]
  \label{eq:boundary_reward}
\end{equation}

where $\tau$ is a regulatory or operational threshold (e.g., 1.8\,m for
fall protection).
The model receives a reward of 1 if it correctly classifies whether the
distance is above or below the threshold, and 0 otherwise.
The combined distance reward is:

\begin{equation}
  R_{\text{spatial}}(\hat{d}, d^*) =
    \lambda_s R_{\text{dist}}(\hat{d}, d^*) +
    \lambda_b R_{\text{boundary}}(\hat{d}, d^*; \tau)
  \label{eq:combined_spatial}
\end{equation}

with $\lambda_s = 0.7$, $\lambda_b = 0.3$ in our experiments.

\subsubsection{Spatial Relationship Reward}

For relationship extraction, we use a set-match reward.
Given the model's predicted relationship set $\hat{S}$ and the
ground-truth set $S^*$ (computed from 3D geometry---bounding box
positions, depth ordering, containment checks):

\begin{equation}
  R_{\text{rel}}(\hat{S}, S^*) =
    \text{F1}(\hat{S}, S^*) =
    \frac{2 |\hat{S} \cap S^*|}{|\hat{S}| + |S^*|}
  \label{eq:rel_reward}
\end{equation}

Each relationship $(e_i, \text{rel}, e_j)$ is treated as a set element.
Precision penalizes hallucinated relationships; recall penalizes
missed ones.

Ground truth is computed programmatically: given detected bounding boxes
and their depth values, we can determine ``A is left of B'' (compare
$x$-coordinates), ``A is above B'' (compare $y$-coordinates with depth
correction for perspective), ``A is in front of B'' (compare depth), and
``A is near B'' (metric distance below threshold).

\subsubsection{Change Detection Reward}

For temporal change detection between frames $f_t$ and $f_{t'}$:

\begin{equation}
  R_{\text{CD}}(\hat{C}, C^*) =
    \text{F1}(\hat{C}, C^*)
    \cdot \mathbf{1}\!\left[\Delta(f_t, f_{t'}) > \epsilon\right]
  \label{eq:cd_reward}
\end{equation}

where $\hat{C}$ is the predicted change set, $C^*$ is the ground-truth
change set (computed by differencing object detections across frames),
and $\Delta$ is a scene-change magnitude measure (optical flow or
detection set difference).
The indicator ensures we only evaluate change detection when genuine
scene change has occurred---otherwise, the trivially correct answer is
``nothing changed.''

\subsubsection{Aggregate Reward}

For a sequence of $T$ frames, the total reward combines all spatial
dimensions:

\begin{equation}
  \mathcal{R}_{\text{total}} =
    \frac{1}{T}\sum_{t=1}^{T}
    \left[\alpha R_{\text{spatial}}(f_t) +
          \beta R_{\text{rel}}(f_t) +
          \gamma R_{\text{CD}}(f_t, f_{t-1})\right]
  \label{eq:total_reward}
\end{equation}

with $\alpha = 0.5$, $\beta = 0.3$, $\gamma = 0.2$ in our experiments.
Distance estimation receives the highest weight because it is the most
directly verifiable and the most operationally consequential.

\subsection{Pipeline Architecture}
\label{sec:pipeline}

\Cref{fig:pipeline} shows the VINNA pipeline.

\begin{figure}[t]
  \centering
  \fbox{\rule{0pt}{4cm}\rule{0.9\linewidth}{0pt}}
  \Description{Block diagram of the VINNA pipeline: egocentric frames
               flow through sensor fusion, spatial VLM analysis, and
               verifiable reward scoring to produce spatial scene graphs.}
  \caption{VINNA pipeline. Egocentric fisheye frames and LiDAR point
           clouds are fused into depth-augmented RGBD frames. A spatial
           VLM emits structured spatial claims (distances, relationships,
           changes, zone classifications). Verifiable reward functions
           score each claim against geometric ground truth. Downstream
           applications (compliance, progress monitoring, resource
           tracking) consume the spatial representation.}
  \label{fig:pipeline}
\end{figure}

\paragraph{Sensor Fusion.}
Each Ironsite bodycam frame (fisheye, $\sim$180$^\circ$ FOV) is paired
with the nearest LiDAR sweep ($\leq 50$\,ms temporal offset).
We project the 3D point cloud into the camera frame using factory
extrinsic calibration, producing a depth-augmented RGBD representation.
Metric distances are read directly from the fused depth channel,
bypassing monocular depth estimation.
Where LiDAR coverage is incomplete (occlusion from equipment, range
limits), we fall back to Depth Anything V2~\cite{depthanything2024}
monocular estimates, flagging these as lower-confidence.

\paragraph{Spatial VLM Analysis.}
The depth-augmented frame is passed to a VLM (Gemini 1.5 Pro in our
prototype) with a structured spatial analysis prompt requesting:
\begin{enumerate}
  \item Scene-level spatial description: overall layout, zone
        identification, activity characterization.
  \item Entity inventory with 3D coordinates: workers, equipment,
        materials, structural elements, each with estimated position
        $(x, y, z)$ in camera-relative meters.
  \item Pairwise spatial relationships: a set of
        $(e_i, \text{rel}, e_j, \text{distance})$ tuples.
  \item For temporal pairs: explicit change list (new/removed/moved
        entities, changed relationships).
  \item Zone-level activity summary: zone type, occupancy count,
        activity level (active/idle/transit).
\end{enumerate}

The VLM output is parsed into a structured JSON spatial scene graph.
This graph is the primary output of VINNA---not a violation flag, but
a \emph{spatial representation} of the construction site at that moment
in time.

\paragraph{Verifiable Reward Scoring.}
Each claim in the spatial scene graph is scored against geometric ground
truth derived from the fused depth data.
Distance claims are scored by \cref{eq:dist_reward}, relationship claims
by \cref{eq:rel_reward}, and change claims by \cref{eq:cd_reward}.
Mismatches between the VLM's spatial claims and the verifiable geometry
constitute the training signal for GRPO optimization.

\paragraph{Downstream Applications.}
The spatial scene graph supports multiple downstream tasks without
task-specific training:
\begin{itemize}
  \item \textbf{Safety compliance}: threshold spatial claims against
        regulatory distances (e.g., worker within 1.8\,m of unguarded
        edge at $>$1.8\,m elevation $\rightarrow$ fall hazard flag).
  \item \textbf{Progress monitoring}: temporal change in structural
        elements across days (new walls, completed floors).
  \item \textbf{Resource tracking}: zone occupancy patterns over time
        (where are workers spending their time?).
  \item \textbf{Equipment routing}: spatial occupancy maps for path
        planning.
\end{itemize}

\subsection{CII Classification as Downstream Task}
\label{sec:cii}

To demonstrate that spatial understanding subsumes violation detection,
we implement CII compliance classification~\cite{cii2020safety} as a
purely downstream consumer of VINNA's spatial output:

\begin{equation}
  \text{CII}(f) =
  \begin{cases}
    \text{NC} & \text{if } \exists\, (e_i, d_{ij}): d_{ij} < \tau_{\text{crit}} \\
    \text{P}  & \text{if } \exists\, (e_i, d_{ij}): \tau_{\text{crit}} \leq d_{ij} < \tau_{\text{warn}} \\
    \text{C}  & \text{otherwise}
  \end{cases}
  \label{eq:cii}
\end{equation}

where $d_{ij}$ is the spatial distance between entity $e_i$ and a
hazard boundary, $\tau_{\text{crit}}$ is the OSHA-specified critical
threshold, and $\tau_{\text{warn}}$ is a warning threshold.
No additional training or compliance-specific reward is needed---the
CII tier falls directly out of the spatial distance estimates.

%% =====================================================================
\section{Experiments}
\label{sec:experiments}

\subsection{Dataset: Ironsite Masonry Footage}
\label{sec:dataset}

Our dataset consists of egocentric footage collected by Ironsite
on an active masonry construction site.
Statistics are summarized in \cref{tab:dataset}.

\begin{table}[t]
  \centering
  \caption{Ironsite dataset statistics.}
  \label{tab:dataset}
  \begin{tabular}{@{}ll@{}}
    \toprule
    \textbf{Property} & \textbf{Value} \\
    \midrule
    Duration          & 21 minutes 14 seconds \\
    Total frames extracted & 638 \\
    Frame rate        & $\sim$0.5 fps (key-frame extracted) \\
    Temporal batches  & 39 \\
    Embedding dim.    & 3072 (Gemini Embedding~2) \\
    Sensor            & Egocentric fisheye bodycam \\
    Depth sensor      & LiDAR point cloud (co-registered) \\
    Task category     & Masonry (block laying, mortar, scaffold) \\
    Human annotations & 25 frames (held-out evaluation only) \\
    \bottomrule
  \end{tabular}
\end{table}

The footage covers workers on scaffolding at multiple heights, masonry
block handling, mortar mixing, and material storage.
The 39 temporal batches represent contiguous segments of related
activity, averaging $\sim$16 frames (32 seconds) per batch.

\subsection{Spatial Analysis Protocol}
\label{sec:protocol}

Each of the 638 frames is processed through the full VINNA pipeline:

\begin{enumerate}
  \item \textbf{Sensor fusion}: fisheye frame + LiDAR $\rightarrow$
        RGBD depth-augmented frame.
  \item \textbf{Spatial VLM analysis}: Gemini 1.5 Pro emits structured
        spatial scene graph (entities, coordinates, relationships,
        zone classification).
  \item \textbf{Embedding}: the full structured JSON output is embedded
        via Gemini Embedding~2 into a 3072-dimensional vector.
  \item \textbf{Reward scoring}: distance claims scored against LiDAR
        depth; relationship claims scored against bbox geometry;
        change claims scored against frame-pair differencing.
  \item \textbf{Downstream CII}: spatial distances thresholded to
        produce P/C/NC classification.
\end{enumerate}

\subsection{Embedding Space Analysis}
\label{sec:embedding}

To validate that VINNA's spatial representations capture meaningful
scene-level structure, we analyze the embedding space of all 638 frames.

\Cref{tab:embedding} reports mean within-cluster cosine similarity
for frames grouped by their downstream CII tier.
If VINNA's spatial representations are informative, frames with similar
spatial conditions should cluster together---even though the embeddings
are computed from \emph{spatial scene graphs}, not compliance labels.

\begin{table}[t]
  \centering
  \caption{Mean within-cluster cosine similarity by spatial condition
           tier (Gemini Embedding~2, 3072-dim). Higher values indicate
           tighter, more coherent spatial representation clusters.}
  \label{tab:embedding}
  \begin{tabular}{@{}lccc@{}}
    \toprule
    \textbf{Spatial Tier} & \textbf{Mean Cos. Sim.} &
    \textbf{Frames} & \textbf{Batches} \\
    \midrule
    Partial (P)         & 0.4246 & 187 & 14 \\
    Conforming (C)      & 0.4127 & 341 & 19 \\
    Non-Conforming (NC) & 0.3765 & 110 &  6 \\
    \bottomrule
  \end{tabular}
\end{table}

The P tier (spatially complex scenes with mixed conditions) shows the
highest intra-cluster coherence, suggesting that ``partial'' spatial
configurations share common structural features (e.g., workers at
moderate distances from hazards, mixed equipment/material layouts).
The NC tier shows the lowest coherence, consistent with the spatial
diversity of hazardous configurations---there are many distinct spatial
arrangements that constitute a non-conforming condition.

\subsection{Spatial Claim Accuracy}
\label{sec:spatial_accuracy}

\Cref{tab:spatial_precision} evaluates the precision and recall of
VINNA's four spatial claim types on the 25-frame annotated subset.

\begin{table}[t]
  \centering
  \caption{Spatial claim accuracy by claim type (25 annotated frames).}
  \label{tab:spatial_precision}
  \begin{tabular}{@{}lcc@{}}
    \toprule
    \textbf{Spatial Claim Type} & \textbf{Precision} & \textbf{Recall} \\
    \midrule
    Entity detection \& localization & 0.94 & 0.91 \\
    Metric distance estimation       & 0.87 & 0.79 \\
    Spatial relationship extraction  & 0.89 & 0.83 \\
    Zone activity classification     & 0.91 & 0.88 \\
    \bottomrule
  \end{tabular}
\end{table}

Entity detection achieves the highest accuracy, benefiting from the
LiDAR point cloud which provides unambiguous evidence of physical
objects.
Metric distance estimation is the weakest spatial capability, primarily
due to LiDAR occlusion: when one or both endpoints of a distance
measurement are occluded in the point cloud, the system falls back to
monocular depth estimates, which are less precise.

\subsection{Downstream: CII Compliance Classification}
\label{sec:cii_results}

As a downstream demonstration, we apply the CII classification rule
(\cref{eq:cii}) to VINNA's spatial distance estimates without any
compliance-specific training.

\begin{table}[t]
  \centering
  \caption{CII compliance classification from spatial representations
           vs.\ expert annotation (25 frames). No compliance-specific
           training---CII tiers are derived entirely from spatial
           distance thresholds.}
  \label{tab:cii_accuracy}
  \begin{tabular}{@{}lccc@{}}
    \toprule
    \multirow{2}{*}{\textbf{Predicted}} &
      \multicolumn{3}{c}{\textbf{Ground Truth}} \\
    \cmidrule(l){2-4}
    & \textbf{C} & \textbf{P} & \textbf{NC} \\
    \midrule
    \textbf{C}  & 11 & 1 & 0 \\
    \textbf{P}  &  1 & 5 & 1 \\
    \textbf{NC} &  0 & 1 & 5 \\
    \bottomrule
  \end{tabular}
\end{table}

Macro-averaged accuracy is 84\% (21/25 correct).
The key observation is that this accuracy is achieved \emph{without any
compliance-specific training}: the system has never seen an OSHA
regulation or a labeled violation.
It achieves compliance classification accuracy purely from spatial
understanding---knowing distances, relationships, and zone configurations
is sufficient to determine whether a scene is conforming.

Most confusions occur at the P/NC boundary, which corresponds to the
spatial ambiguity zone where distances are near regulatory thresholds.
Improved LiDAR depth calibration would directly improve this boundary
discrimination.

\subsection{Representative Spatial Analyses}
\label{sec:events}

\Cref{tab:events} presents five representative spatial analyses,
showing the structured spatial outputs alongside derived downstream
information.

\begin{table*}[t]
  \centering
  \caption{Representative spatial analyses from Ironsite masonry
           footage. Columns show spatial outputs (distances,
           relationships) and derived downstream CII classification.}
  \label{tab:events}
  \small
  \begin{tabular}{@{}cllccc@{}}
    \toprule
    \textbf{Frame} & \textbf{Spatial Description} &
    \textbf{Key Distance} & \textbf{Relationships} &
    \textbf{CII} & $R_{\text{spatial}}$ \\
    \midrule
    042 & Worker on scaffold, edge 0.7\,m, elev.\ 4.2\,m &
          worker--edge: 0.7\,m & \texttt{on(scaffold), near(edge)} &
          NC & $0.18$ \\
    \addlinespace
    117 & Block laying at ground, nearest hazard 8.3\,m &
          worker--hazard: 8.3\,m & \texttt{on(ground), using(block)} &
          C  & $0.94$ \\
    \addlinespace
    231 & Mortar mixing, equipment 1.4\,m, no eye protection &
          worker--equip: 1.4\,m & \texttt{near(mixer), in(work\_zone)} &
          P  & $0.69$ \\
    \addlinespace
    389 & Worker near roof edge, tethered, dist.\ 2.1\,m &
          worker--edge: 2.1\,m & \texttt{near(edge), tethered\_to(line)} &
          C  & $0.89$ \\
    \addlinespace
    512 & Materials stacked, worker 0.9\,m from 3.8\,m drop &
          worker--drop: 0.9\,m & \texttt{near(edge), near(materials)} &
          NC & $0.24$ \\
    \bottomrule
  \end{tabular}
\end{table*}

Note how the spatial representation is richer than a violation flag.
Frame 389 illustrates this: a violation detector would flag or clear the
scene, but the spatial representation captures that the worker is near
the edge (2.1\,m) \emph{and} tethered, the edge is at roof height, and
the worker's activity context is ``roofing.''
This spatial richness enables multiple downstream uses of the same
frame analysis.

\subsection{Temporal Batch Structure}
\label{sec:temporal}

The 39 temporal batches provide a natural structure for evaluating
spatial consistency over time.
Within each batch (contiguous frames of related activity), we expect
spatial relationships to be relatively stable, with changes occurring
primarily through worker movement and material manipulation.

We measure intra-batch vs.\ inter-batch embedding similarity:

\begin{table}[t]
  \centering
  \caption{Embedding similarity: intra-batch vs.\ inter-batch.}
  \label{tab:temporal}
  \begin{tabular}{@{}lc@{}}
    \toprule
    \textbf{Comparison} & \textbf{Mean Cosine Sim.} \\
    \midrule
    Intra-batch (same activity) & 0.5832 \\
    Inter-batch (different activity) & 0.3214 \\
    Adjacent batches & 0.4107 \\
    \bottomrule
  \end{tabular}
\end{table}

The 0.26 gap between intra-batch and inter-batch similarity confirms
that the spatial representations capture activity-level scene structure.
Adjacent batches show intermediate similarity, consistent with gradual
spatial transitions between activities (e.g., the scene does not change
abruptly when a worker transitions from block-laying to mortar mixing).

%% =====================================================================
\section{Discussion}
\label{sec:discussion}

\subsection{Spatial Intelligence vs.\ Violation Detection}

The core argument of this paper is that construction AI should target
\emph{spatial intelligence}, not narrow classification.
Our results provide preliminary evidence for this claim:

\begin{itemize}
  \item CII compliance classification emerges from spatial distance
        estimates without task-specific training (84\% accuracy).
  \item The spatial representation captures scene structure
        (embedding cluster separation) that a binary violation flag
        would destroy.
  \item Temporal batch coherence demonstrates that spatial
        representations track activity-level scene evolution.
\end{itemize}

The practical implication is architectural: build the spatial
reasoning layer first, then derive any number of downstream
applications from it.
Do not build a violation detector and then try to reverse-engineer
spatial understanding from its classification boundary.

\subsection{Differentiation from Safe-Construct}

Safe-Construct~\cite{safeconstruct2025} is the most directly relevant
prior work.
The key difference is not methodological but \emph{ontological}:
Safe-Construct treats violation detection as the output.
VINNA treats spatial understanding as the output and violation detection
as a derived quantity.

This matters practically because:
\begin{enumerate}
  \item Spatial representations are \emph{reusable}. The same spatial
        scene graph that detects a fall hazard also tracks progress,
        monitors equipment utilization, and supports logistics planning.
  \item Spatial rewards are \emph{more verifiable} than violation
        labels. ``The distance is 2.3\,m'' is checkable against a depth
        sensor. ``This is a violation'' requires regulatory
        interpretation.
  \item Spatial understanding \emph{degrades gracefully}. A distance
        estimate can be approximately right; a violation flag is right
        or wrong.
\end{enumerate}

\subsection{Verifiable Spatial Rewards vs.\ Human Labels}

The practical advantage of verifiable spatial rewards is
\emph{scalability without annotation}.
A human annotation pipeline for spatial relationships in 638 frames
would require expert time for each entity pair in each frame---easily
hundreds of hours.
Our verifiable reward pipeline computes ground truth from sensor data
in under 4 minutes (VLM inference is the bottleneck).

A subtler advantage is \emph{reward granularity}.
Human annotators can provide binary labels (correct/incorrect) or at
best ordinal ratings.
SNRA distance rewards provide a continuous signal that tells the model
not just that its estimate was wrong, but \emph{how wrong} and \emph{in
which direction}---exactly the gradient information GRPO needs.

\subsection{Limitations}

\paragraph{No trained RL model.}
The most significant limitation is that we have not yet executed the
GRPO training loop.
Our reward functions are designed and scored, but the reinforcement
learning step---actually updating a VLM's weights using verifiable
spatial rewards---remains future work.
The results reported here are from the Gemini 1.5 Pro VLM in a
zero-shot structured prompting configuration, not a spatially fine-tuned
model.

\paragraph{Small dataset.}
638 frames from a single masonry site is insufficient for claims about
generalization.
Different construction trades (steel erection, roofing, concrete
formwork, electrical) have different spatial configurations.
We expect the spatial reasoning capabilities to transfer (distances and
relationships are domain-agnostic), but the zone activity classifier
will require trade-specific tuning.

\paragraph{LiDAR occlusion.}
The primary failure mode for distance estimation is LiDAR occlusion:
when an entity is not visible in the point cloud (behind equipment,
outside LiDAR range), the system falls back to monocular depth
estimates, which are less reliable.
Multi-camera fusion or floor-plane priors would mitigate this.

\paragraph{Fisheye distortion.}
The 180$^\circ$ fisheye lens produces significant radial distortion
that the VLM must implicitly correct.
We observe that spatial relationship accuracy degrades for entities
near the frame periphery, where distortion is most severe.

\paragraph{Evaluation scale.}
Our annotated evaluation set (25 frames) is small.
We report it honestly as a directional signal, not a statistically
robust benchmark.

\subsection{Future Work}

\paragraph{GRPO training.}
The immediate next step is executing the GRPO training loop using
our designed reward functions.
Following Faithful GRPO~\cite{faithfulgrpo2026}, we plan to fine-tune
Qwen2.5-VL-7B with LoRA (rank 16--32) using the spatial distance
reward as the primary training signal and relationship/change rewards
as auxiliary signals.

\paragraph{Temporal modeling.}
Our current analysis is frame-level.
Extending to explicit temporal modeling---tracking entities across
frames, maintaining persistent spatial maps, detecting trends in zone
activity---would enable richer construction intelligence.

\paragraph{Multi-site evaluation.}
Generalizing beyond masonry to other construction trades, and from a
single site to multiple sites, is essential for validating the claim
that spatial intelligence is a general foundation.

\paragraph{Real-time spatial inference.}
For deployment, the VLM inference bottleneck must be addressed.
A cascaded architecture---lightweight spatial front-end (depth-based
distance estimation, YOLO entity detection) for high-frequency frames,
full VLM analysis triggered on spatial anomalies---mirrors the
wake-word architecture used in speech systems.

%% =====================================================================
\section{Conclusion}
\label{sec:conclusion}

Construction AI has been building the house from the roof down:
jumping to narrow classification tasks (PPE detection, violation
flagging) without first laying the foundation of spatial understanding.
VINNA inverts this, treating spatial intelligence---metric distance
estimation, spatial relationship extraction, temporal change detection,
zone activity analysis---as the primary capability and classification
as a downstream application.

On 638 frames of real Ironsite egocentric footage from an active
masonry site, VINNA's spatial representations produce meaningful
embedding-space structure across 39 temporal batches.
CII compliance classification emerges from spatial distance estimates
at 84\% accuracy without compliance-specific training, demonstrating
that spatial understanding subsumes violation detection.

Our verifiable reward design---distance checked against LiDAR,
relationships validated by 3D geometry, changes scored by frame
differencing---eliminates the need for human annotation and enables
self-supervised GRPO training.
This is a hackathon prototype with real data, not a finished system.
The spatial intelligence gap in construction AI is real, and VINNA is
a first step toward closing it.

%% =====================================================================
\begin{acks}
The authors thank Ironsite for access to the egocentric bodycam and
LiDAR dataset, and the HackTech 2025 organizers at Caltech for hosting
this work.
\end{acks}

%% =====================================================================
\begin{thebibliography}{99}

\bibitem{spatialvlm2024}
Zhiyuan Chen, Jiahao Li, and Kilian Q. Weinberger.
\newblock {SpatialVLM}: Endowing vision--language models with spatial reasoning
  capabilities.
\newblock In \emph{Proceedings of the IEEE/CVF Conference on Computer Vision
  and Pattern Recognition (CVPR)}, 2024.

\bibitem{spatialrgpt2024}
Rongyao Fang, Yingtian Tang, and Hao Chen.
\newblock {SpatialRGPT}: Grounded spatial reasoning in vision--language models.
\newblock In \emph{Advances in Neural Information Processing Systems
  (NeurIPS)}, 2024.

\bibitem{faithfulgrpo2026}
Sai~Srinivas Kancheti, Aditya Kanade, Rohit Sinha, Vineeth~N Balasubramanian, and Tanuja Ganu.
\newblock Faithful {GRPO}: Improving visual spatial reasoning in multimodal
  language models via constrained policy optimization.
\newblock \emph{arXiv:2604.08476}, April 2026.

\bibitem{smoothoperator2025}
Siwen Jiao, Tianxiong Lv, Kangan Qian, Chenxu Zhao, Xiuyuan Zhu, Tianlun Li,
  Xiaolong Cheng, Jinyu Li, Zhihao Liao, and Yang Cai.
\newblock Smooth operator: Smooth verifiable reward activates spatial reasoning
  ability of vision-language models.
\newblock \emph{arXiv:2601.07695}, January 2026.

\bibitem{ego4d2022}
Kristen Grauman et~al.
\newblock Ego4{D}: Around the world in 3,000 hours of egocentric video.
\newblock In \emph{Proceedings of the IEEE/CVF Conference on Computer Vision
  and Pattern Recognition (CVPR)}, 2022.

\bibitem{safeconstruct2025}
Aviral Chharia, Tianyu Ren, Tomotake Furuhata, and Kenji Shimada.
\newblock {Safe-Construct}: Redefining construction safety violation
  recognition as 3D multi-view engagement task.
\newblock In \emph{CVPR Workshops}, 2025.
\newblock arXiv:2504.10880.

\bibitem{vigorl2025}
Gabriel Sarch, Snigdha Saha, Naitik Khandelwal, Ayush Jain, Michael~J. Tarr,
  Aviral Kumar, and Katerina Fragkiadaki.
\newblock Grounded reinforcement learning for visual reasoning.
\newblock \emph{arXiv:2505.23678}, 2025.

\bibitem{depthanything2024}
Lihe Yang et~al.
\newblock Depth Anything V2.
\newblock In \emph{Advances in Neural Information Processing Systems
  (NeurIPS)}, 2024.

\bibitem{golparvar2011monitoring}
Mani Golparvar-Fard, Feniosky Pe\~{n}a-Mora, and Silvio Savarese.
\newblock Integrated sequential as-built and as-planned representation with
  backward and forward differencing approach.
\newblock \emph{Journal of Computing in Civil Engineering}, 25(2):98--116,
  2011.

\bibitem{yang2018worker}
Jingjing Yang, Pingbo Tang, and Changbum Ahn.
\newblock Construction worker detection in video frames for project
  performance monitoring.
\newblock \emph{Automation in Construction}, 92:2--16, 2018.

\bibitem{nath2020ppe}
Nipun~D. Nath, Amir~H. Behzadan, and Stephanie~G. Paal.
\newblock Deep learning for site safety: Real-time detection of personal
  protective equipment.
\newblock \emph{Automation in Construction}, 112:103085, 2020.

\bibitem{christiano2017deep}
Paul~F. Christiano, Jan Leike, Tom Brown, Miljan Martic, Shane Legg, and
  Dario Amodei.
\newblock Deep reinforcement learning from human preferences.
\newblock In \emph{Advances in Neural Information Processing Systems
  (NeurIPS)}, 2017.

\bibitem{cii2020safety}
{Construction Industry Institute}.
\newblock Safety performance measurement system, 3rd edition.
\newblock Research Report {CII-IR174-3}, University of Texas at Austin, 2020.

\end{thebibliography}

\end{document}
